\documentclass[main.tex]{subfiles}
\begin{document}

\label{subsec:data}

I combine five data sources to measure consumer reward sensitivity, how card acceptance affects sales, and multi-homing rates.
Appendix \ref{sec:appendix-data} provides data construction details.

\vspace{0.05in}

\noindent \textbf{Aggregate Prices and Shares (2019)}: I use the Nilson Report's aggregate payment volumes and merchant fees in 2019.
I collect average rewards data from the annual reports of the $6$ largest banks,
% CONSTANT: data collection method, count of banks scraped
whose cards cover 80\% of credit card purchase volumes.
% CONSTANT: top banks cover ~80% of credit card purchase volumes, from Nilson Report
These data pin down baseline market shares and fee levels in the model.
Figure \ref{fig:aggregate-shares} shows payment volumes, merchant fees, and rewards.
Visa, MC, and AmEx process $90\%$ of all card payments.
% CONSTANT: Visa/MC/AmEx process ~90% of card payments, from Nilson Report
All three major credit card networks charge merchant fees of around $2.25\%$,
% CONSTANT: average credit card MDR ~2.25%, from Nilson Report
whereas debit networks charge around $0.7\%$ due to the Durbin Amendment.
% CONSTANT: average debit MDR ~0.7%, from Nilson Report
Credit cards pay rewards around 1.4\%, whereas aggregate debit rewards are effectively zero.
% CONSTANT: average credit card reward rate from bank annual reports, external data

\begin{figure}
    \small
    \caption{Aggregate payment volumes, merchant fees, and rewards in 2019}
    \label{fig:aggregate-shares}
    \begin{center}
        \includegraphics[width=2.5in]{../output/graphs/volume_plot.pdf}
        \includegraphics[width=2.5in]{../output/graphs/price_plot.pdf}
    \end{center}
    \tablenote{The left chart shows payment volumes measured in trillions from \textcite{Nilson2020a,Nilson2020b}.
Visa and MC operate credit and debit card networks, whereas AmEx primarily offers credit and charge cards.
Discover is much smaller than the other three networks.
The right chart shows merchant fees from \textcite{Nilson2020} and rewards from banks' annual financial reports.
Aggregate debit rewards fell to effectively zero after the Durbin Amendment, though some small, exempt issuers continued to offer them \parencite{Hayashi2012}.
The cost of cash is from \textcite{Felt.etal2020}.}

\end{figure}

\vspace{0.05in}

\noindent \textbf{Issuer Payment Volumes (2007--2014).} I construct an annual panel of 39 issuers from the Nilson Report covering 2007--2014 to study how the Durbin Amendment affected payment volumes.
The issuers include both banks and large credit unions.
% HARDCODED[scalar: nilson_num] current=39
I merge this data with NCUA call reports for credit unions and Y-9C reports for bank holding companies to get issuer-level financial data.

\vspace{0.05in}

\noindent \textbf{Homescan (2013--2023).} The NielsenIQ Homescan panel tracks the payment decisions of around $100,000$ households at large consumer packaged goods stores from 2013--2023.
% CONSTANT: Homescan panel size, from Nielsen data documentation
I use this panel to measure how consumers split spending across the cards in their wallets and how card acceptance affects sales.
Appendix Table~\ref{tab:nielsen-summary} reports summary statistics.
Appendix Table~\ref{tab:nielsen-compare} compares Homescan payment shares to aggregate shares.
The Homescan sample is broadly representative except for AmEx, which is underrepresented at 4\% versus 10\% aggregate share. This reflects Homescan's exclusion of sectors like travel where AmEx dominates.
% HARDCODED[scalar: amex_uncond_share] current=4
% CONSTANT: 10% is AmEx aggregate Nilson share, external data source
For this reason, AmEx-specific calibration moments (merchant fees, rewards) are sourced externally rather than from Homescan.

\vspace{0.05in}

\noindent \textbf{Consumer Payment Surveys.} I combine the 2017--2019 waves of the Atlanta Federal Reserve's Diary of Consumer Payment Choice (DCPC) and its companion Survey (SCPC) to build a transaction-level dataset over three-day windows \parencite{Greene.Stavins2021}.
The rich demographic detail in these surveys identifies income gradients in payment preferences.
The diary and survey share a common panel of respondents. 
Whereas Homescan oversamples large retailers, the DCPC captures transactions at small stores, allowing me to better estimate overall card acceptance rates.
Table \ref{tab:dcpc-summary} shows summary statistics on consumers' payment preferences.
Cards are accepted for around 95\% of transactions.\footnote{
    Online Appendix Table~\ref{tab:yelp-dcpc-comparison} shows that diary-based acceptance rates closely match those from Yelp business listings.
}
% HARDCODED[scalar: share_card_cons] current=95
Debit cards are the most popular payment method.
Credit card users have higher incomes than cash or debit card users, and around one-quarter of transactions are online.
% Cash-only merchants account for the remaining 5\% of transaction volume, concentrated in small service providers and informal sellers. Around 25\% of consumers prefer cash as their primary payment method (Table~\ref{tab:dcpc-summary}). The model treats online and in-person transactions symmetrically; because virtually all online merchants accept cards, the 25\% online share reinforces the high baseline acceptance rate.
I also conduct a second-choice survey in 2024 to estimate substitution patterns across payment methods \parencite{Berry.etal2004}.

\begin{table}
    \small
        \caption{Summary statistics by consumer type in the Diary of Consumer Payment Choice (DCPC)}
        \label{tab:dcpc-summary}
        \footnotesize
        \begin{center}
        \input{../output/tables/cc_balances_by_pref.tex}
    \end{center}
    \tablenote{Consumers are split into three groups by preferred non-bill payment instrument.
Share reports the fraction of the sample in each column.
Card acceptance is the expenditure share at merchants that accept cards.
All other variables report consumer-level averages.
Online transaction reports the share of transactions done online.}

\end{table}

\end{document}